\begin{exercise}[Asian call versus European call]
\textbf{Restatement.}
For the Asian payoff
\[
  G=\pos{\frac1T\int_0^T S_t\,\dd t-K},
\]
show that its time-\(0\) price is no larger than the European call price with
the same strike and maturity.

\begin{solution}
By risk-neutral pricing, \polyref{Eq.~(3.9)},
\[
  p_0^{\mathrm{Asian}}
    =\EQ\left[
       \e^{-rT}\pos{\frac1T\int_0^T S_t\,\dd t-K}
     \right].
\]
The function \(y\mapsto\pos{y-K}\) is convex, so Jensen's inequality gives the
pathwise bound
\[
  \pos{\frac1T\int_0^T S_t\,\dd t-K}
  \le
  \frac1T\int_0^T\pos{S_t-K}\,\dd t.
\]
Thus
\[
  p_0^{\mathrm{Asian}}
  \le
  \frac1T\int_0^T\e^{-rT}\EQ[\pos{S_t-K}]\,\dd t.
\]
Since \(r\ge0\) in the BMS chapter and \(S\) is a \(\Q\)-submartingale, the
increasing convex function \(x\mapsto\pos{x-K}\) also yields a submartingale:
for \(t\le T\),
\[
  \EQ[\pos{S_T-K}\,\mid\,\F_t]
    \ge \pos{\EQ[S_T\,\mid\,\F_t]-K}
    =\pos{\e^{r(T-t)}S_t-K}
    \ge \pos{S_t-K}.
\]
Taking expectations,
\[
  \EQ[\pos{S_t-K}]\le \EQ[\pos{S_T-K}].
\]
Therefore
\[
  p_0^{\mathrm{Asian}}
    \le \e^{-rT}\EQ[\pos{S_T-K}]
    =p_0^{\mathrm{European}}.
\]
\end{solution}
\end{exercise}
